\documentclass[english]{tudscrartcl}
\usepackage{babel}
\usepackage[babel]{microtype}

\usepackage[utf8]{inputenc}
\usepackage[T1]{fontenc}
\usepackage{clrscode}
\usepackage{xcolor}
\usepackage[pdftex]{graphicx}
\usepackage{xurl}

\begin{document}

	\addtokomafont{subject}{\sffamily}
	\title{Exposé: "Realization of a smart meter prototype based on hardware specifications"}
	\author{Jakob Fregin}
	\date{December 2025}

  	\maketitle
  	\pagenumbering{arabic}
	\section{Introduction}
	\normalsize
	% Intro to SM in SGs
	The rise of smart grids (SG) and consequently the increasing use of smart meters (SM) allows resource efficient balancing between consumption and production inside the grid \cite{en15196984}.
	SMs constantly monitor and report their sensor values, enabling a on-demand, more stable and resilient grid \cite{en15196984}.
	These SMs that formally contained primitive 8-bit MCUs, now feature 32-bit Systems-on-Chip (SoC) with multiple cores, various communication standards and other dedicated hardware \cite{ARTICLE:3}.
	In their basic forms, SMs feature outage detection, remote management, trouble shooting and automatic billing \cite{7379940}.
	Their real-time nature and capability for high sample rates enable fine-grained data analysis, e.g. for outage prediction \cite{8646486}
	or value added services, e.g. dynamic pricing, allowing users to pay the best price available in each measurement interval.
	
	\section{Problem}
	\normalsize
	% Anonymity motivation
	Despite growing capabilities that enhance their versatility and functionality,
	SMs are still exposed to the Internet and the same level of security issues as other usually much more powerful network attached devices. 
	At the same time, SMs pose an imminent threat to peoples privacy, since they collect a significant amount of personal data \cite{7959167},
	due to their high sample rates and applicability for multiple resources, e.g. power, natural gas and water supply.
	For instance, this data infers the presence of people at home or their economic stand based on their resource consumption \cite{7959167}
	or can be further analyzed to extract individual appliance usage \cite{20101011}, exposing users habits and routines \cite{20101011}.
	\\ % SM systems state of law - Germany
	In the current state, as of November 2025 in Germany, smart metering is tightly regulated,
	as SMs are designated to measure the resource at most every 15min (§55 Abs. 1 MsbG) \cite{Msbg55}.
	After aquisition, the data is encrypted as well as signed, ensuring a confidential
	and tamperproof communication between the SM and SMGW (§52 Abs. 1 MsbG) \cite{Msbg52}.
	However, the SMGW can access unencrypted personalized data from multiple SMs.
	Which, if compromised, would expose detailed consumption data of all connected households.
	Only during aggregation anonymization is applied as far as the use case allows (§52 Abs. 3 MsbG) \cite{Msbg52}.
	Apart from that, the granularity of 15min in combination with dynamic pricing allows for detailed billing.
	This detail of billing is another privacy concern, since it potentially exposes the complete consumption behavior of a household,
	while being recognised as valid data monitoring case by law (§50 Abs. 2 10. MsbG) \cite{Msbg50}.
	\\ % Anonymity proposed solution
	We propose that these privacy concerns could be solved by introducing anonymization directly on the SM,
	keeping the data personalized on the customers' premises only. 
	An exemplary concept to realize these privacy-preserving measurements, could include high frequency measurements, encryption, dynamic pricing,
	anonymization amongst peers and evidence for the current price by a zero knowledge proof.
	This concept would eliminate the need for sample rate limits or detailed billing, since privacy and validity are provided by cryptography.
	Any privacy-preserving solutions would inherently result in higher computational effort for the SM as well as more network traffic,
	for which the current SMs might lack in hardware.
	\\ % Energy concerns
	Another concern is energy efficiency \cite{RASHEDMOHASSEL2014473}.
	Especially when creating decetrialized solutions that shall be applicable in SG where every household has at least one SM.
	If SMs were to do extensive cryptography on-device, communicating in elevated frequency with neighbors,
	this would result in longer active states and therefor higher power consumption.
	Regarding SMs that are not directly connected to the power grid this becomes even more crucial,
	since there additional maintenance, i.e. swapping a battery, would arise.
	Therefor, an energy efficient approach would not only be more environmentally friendly,
	but also in the interest of resource providers and customers.
	For this, it is important to quantify resources and consumption of the SMs.
	\\ % Final Problem / Goal Spec
	In the upcoming work, we want to research the harware capabilities of SMs,
	in order to explore the feasibility for extended on-device computation.
	Therefore we want to create a representative SM hardware emulation,
	benchmark it with computational operations and analyze the results,
	to provide an outlook on the viability of the SM being a stand-alone, secure, privacy concerned device,
	sharing only aggregated information as needed for billing or grid management.
	
	\section{Related Work}
	% OSM
	There already is an Open Smart Meter project, which is an open source hardware smart meter design \cite{OSM-WebSite, OSM-GitHub},
	which can be used for SM research, experiemnts or deployment.
	But since it is a HW SM, we still see a benefit in emulating the HW, to avoid the HW build step and setup.
	\\ % MCU emulation
	Furthermore there are various MCU emulators available, like e.g. QEMU \cite{QEMU-WebSite}, a functional instruction-set simulator(ISS) supporting several architectures,
	which already have been used for emulating IoT devices \cite{66920, 10.1007/978-3-031-47457-6_5}.
	For a more detailed hardware representation, there are system emulators like e.g. gem5 \cite{gem5-WebSite} or SystemC \cite{1193229},
	which allow detailed models of CPUs, memory and peripherals. We consider those as options or references for our emulation,
	while our focus will be on algorithm emulation including acurrate benchmarking, instead of real-time or explicit hardware simulation.
  	\\ % IoT Benchmarking
  	There also have been proposals for IoT device benchmarks \cite{6945583, 8802949}, like e.g. BenchIoT \cite{8809492}.
	This benchmark is built to run on devices, comparable to SM MCUs, e.g. of the ARM Cortex M series.
	Additionally, there are benchmarks for embedded devices in general, like e.g. EEMBC \cite{EEMBC-WebSite},
	which provides computational workloads, such as CoreMark.
	We want to use those benchmarks as references or even directly on our emulated SM.
	\\ % Conclusion
	Those Tools offer a starting point and reference for our emulation and benchmarking approach,
	while some of their points, especially quite detailed hardware considerations, go out of our scope.

	\section{Our Approach}
	\normalsize
  	% Model
	We want to analyse the current hardware specifications of residential SMs,
	in order to model and build a representative prototype as a virtual machine or in another emulated fashion.
	Our work shall be based on publicly available information and literature,
	as well as gathered specifications by manufactures or government regulatory institutes, if they are given on request.
	Depending on the variety of harware specifications, we want to choose a representative set or devide the model into e.g. high-end and low-end.
	Based on this, the chosen parameters shall represent a state of the art residential SM as closely as possible,
	considering parameters like cpu cores and instructions per second, chache size, memory size, storage size, networking capabilites and energy consumption. 
	\\ % Validation
  	In the following step, we want to validate the emulated SM as accurately as possible.
	First we will try to acquire actual HW benchmark data from vendors or government regulatory institutes,
	to rerun those on the emulation as well and then cross check results.
	If benchmarking data is not available as ground truth, we will search literature for similar hardware,
	that may have been benchmarked with published results, to then again, check if those are reproduced by the emulation.
	As a last resort, we would ensure that our emulation does implement the representative parameters chosen in the first step and
	proceed with the analysis of this representative emulator.
	\\ % Benchmarking - TODO
	Given a working emulation, we want to benchmark computational and cryptographic operations to analyze the SMs' capabilites.
	First we want to run established benchmarks, e.g. EEMBCs CoreMark, on the emulation, to present general performance data.
	Then, we choose relevant cryptographic algorithms for encryption, signing and optionally anonymization,
	that are either already in use in SM systems, e.g. AES-128 as stated by the Open metering System Group \cite{OMS-Spec},
	or fit our mentioned goals in terms of anonymization. Those algorithms will then be run on the emulated SM,
	by either compiling their naive implementations for best repeatability or using established benchmarks.
	This data will provide an outlook on SMs capabilities, especially the estimated upper bound to the sampling frequency,
	at which the SM can still encrypt and sign the data in real-time.
	Additionally we want to analyze the network traffic generated by SMs,
	when sending encrypted and signed measurement reports and optionally the overhead introduced by anonymization.
	Since networking depends heavily on the used link, as well as its environment,
	we want to take a theoretical approach to determine maximum data rates,
	by analyzing the needed traffic for reporting measurements.
	\\ % Optional Benchmarks
  	If the work is finished early, we want to run further benchmarks, especially anonymization algorithms, e.g. deZent, on the emulator,
	simulating other nodes and keeping network restictions, to get a first impression whether anonymization on SMs is even viable.
	\\ % Energy
  	This benchmark data would also allow for energy consumption estimations, using active time or energy per cycle characteristics,
	providing an outlook on potential energy overhead, caused by future extensions.
	
	\section{Risks}
	\normalsize
  	% Time
  	As risks, our time constrain of 20 weeks could significantly decrease the extent of work.
	If thats the case, we consider to cut mainly in variety of benchmarks and pilot data created based on the model,
	ensuring that parameters are chosen and the model is validated properly.
	\\ % Setbacks e.g. benchmarking
	Additionally setbacks in task completion as specified in the timeline, e.g. running a certain benchmark implemtation on the emulator, could arise.
	We plan to always progress in parallel on multiple tasks, e.g. start data aquisition early on benchmarks that already work, to avoid blocking the whole work on a single issue.
	\\ % HW limits
	In general, we consider that, todays hardware might be neither performant nor effiecient enough for a full on device anonymization.
	While this being the point of creating a emulated HW model, it could be that is not that much spare computation available on current devices.
	This would still leave the model and emulator as a quick and easy tool to check whether other concepts are feasible, without extensive prototyping.
  	\\ % Real world implementation
	Another concern is real world implementation.
	Our goals are primarily customer-focused, since we aim to provide users with greater privacy and control over their data.
	However, this approach requires additional power and extended software support, increasing service maintenance costs.
	For these reasons, we think that, adopting this type of system may not be attractive to service providers.
	
	\section{Timeline}
	\normalsize
	Since this work, as described, already breaks down into the following steps we allocate the maximum reserve of 20 weeks preliminarily as follows (listed as Work / Writing):
	\begin{enumerate}
		\item Collection and analysis of SM hardware data / Related Work + Background (3 weeks)
		\item Modelling and building of a representative simulation / Related Work + Background + Evaluation(Log) (3 weeks)
		\item Choice, analysis and implementation of cyptographic algorithms in SM systems / Model + Methodology + Evaluation(Log) (3 weeks)
		\item Benchmarking and data collection using the emulation / Evaluation (3 weeks)
		\item Conclusion: data analysis, forming results / Plots + Discussion + Limitations + Conclusion (3 weeks)
		\item Paper finalization: reworks and polishing / Introduction + Future Work + Abstract (5 weeks)
	\end{enumerate}
	
	\bibliographystyle{plain}
	\bibliography{references}

\end{document}
